%! AP Statistics — Unit 5, Lesson 5.1 — Exit Ticket (Answer Key)
\documentclass[10pt]{article}
\usepackage{apstats-article}
\usepackage{apstats-key}   % pulls in -boxes; do NOT also load -boxes
\newcommand{\TallMath}[1]{$\displaystyle #1\rule[-1.4em]{0pt}{3.2em}$}

\begin{document}

\pageheader{Unit 5, Lesson 5.1}{Exit Ticket --- Answer Key}
\namedateperiod

\vspace{0.15in}

\begin{enumerate}
  \item A nationwide survey of 1,000 adults finds that 62\% support a new education
        policy. A local survey of 50 adults from the same population finds 55\%.
        Explain why these results are \emph{not} necessarily contradictory.

        \ansline{Both are statistics from random samples of the same population. Different random samples will produce different sample proportions just by chance (sampling variability). The 7-percentage-point difference is consistent with what we expect when comparing a large sample to a small sample from the same population.}

  \item Identify each value below as a statistic or a parameter.
        \begin{enumerate}[label=\alph*.]
          \item The true proportion of all U.S. adults who support the policy.
                \hfill \ans{parameter}
          \item The 62\% reported by the nationwide survey.
                \hfill \ans{statistic}
        \end{enumerate}

  \item In your own words, what is a \textbf{sampling distribution}?

        \ansline{A sampling distribution is the distribution of all values a statistic can take when computed from all possible samples of a given size from the same population.}
\end{enumerate}

\end{document}
